\section{从注意力到 Transformer 与大语言模型}

\subsection{注意力的历史演进}

注意力机制从2014年的发明到成为现代 AI 的基石，经历了快速而深远的发展：

\begin{table}[H]
\centering
\begin{tabular}{lll}
\toprule
\textbf{年份} & \textbf{里程碑} & \textbf{关键人物} \\
\midrule
2014 & 加法注意力（NMT） & Bahdanau, Cho, Bengio \\
2015 & 乘法/双线性注意力 & Luong, Pham, Manning \\
2017 & Transformer（自注意力） & Vaswani et al. \\
2018 & BERT（双向 Transformer） & Devlin et al. \\
2018--今 & GPT 系列（大语言模型） & OpenAI \\
\bottomrule
\end{tabular}
\caption{注意力机制的发展时间线}
\end{table}

\begin{importantbox}{注意力是深度学习新时代的开端}
Manning 教授指出，此前神经网络的所有核心组件（前馈网络、RNN、LSTM、CNN）都是在2000年之前发明的，深度学习革命的前半段主要是``等待数据和算力赶上''。而注意力是2014年之后\textbf{真正的新发明}，它开启了 Transformer、BERT、GPT 等一系列变革性工作。
\end{importantbox}

\subsection{从编码器-解码器注意力到自注意力}

在 Seq2Seq 中，注意力用于解码器``回头查看''编码器。而 Transformer 的核心创新是\textbf{自注意力}（Self-Attention）——让序列中的每个位置都可以关注同一序列中的所有其他位置。

\begin{center}
\begin{tikzpicture}[
    node distance=1.5cm,
    every node/.style={draw, rounded corners, minimum width=4.5cm, minimum height=0.8cm, font=\small, align=center},
    arrow/.style={-{Stealth[length=3mm]}, thick}
]
\node (enc-dec) {编码器-解码器注意力\\（Seq2Seq, 2014）};
\node (self) [below=of enc-dec] {自注意力\\（Transformer, 2017）};
\node (multi) [below=of self] {多头注意力\\（并行多组低维注意力）};
\node (llm) [below=of multi] {大语言模型\\（GPT, BERT 等）};

\draw[arrow] (enc-dec) -- (self);
\draw[arrow] (self) -- (multi);
\draw[arrow] (multi) -- (llm);

\node[draw=none, right=0.5cm of enc-dec, text width=5cm, font=\scriptsize, align=left] {解码器查看编码器\\Query=解码器, Values=编码器};
\node[draw=none, right=0.5cm of self, text width=5cm, font=\scriptsize, align=left] {每个位置查看所有位置\\Query=Key=Value=同一序列};
\node[draw=none, right=0.5cm of multi, text width=5cm, font=\scriptsize, align=left] {降秩投影 + 并行多组\\= 降秩乘法注意力的推广};
\node[draw=none, right=0.5cm of llm, text width=5cm, font=\scriptsize, align=left] {堆叠多层 Transformer\\预训练 + 微调/提示};
\end{tikzpicture}
\end{center}

Manning 教授预告，下一讲（Lecture 8）将详细介绍 Transformer 架构，包括自注意力、多头注意力、位置编码等核心概念。

\subsection{大语言模型简介}

课程后半段还涉及了大语言模型（LLM）时代对 NLP 研究和实践的影响：

\begin{itemize}
    \item \textbf{BERT}：基于 Transformer 编码器的双向预训练模型，通过微调适应各种下游任务
    \item \textbf{GPT 系列}：基于 Transformer 解码器的自回归语言模型，展现出强大的零样本和少样本学习能力
    \item \textbf{上下文学习}（In-Context Learning）：大模型可以通过提示（prompt）中的少量示例学习新任务，无需更新参数
\end{itemize}

\begin{knowledgebox}{LLM 时代的研究范式转变}
Manning 教授提到，LLM 的出现改变了研究方式：许多项目不再需要从头训练模型，而是通过 API 调用大语言模型、进行上下文学习或微调来完成任务。这使得更多研究精力可以投入到理解模型行为、评估能力边界、探索新的交互范式等方面。同时，较小的开源模型（如7B参数级别）也为预算有限的研究者提供了可能。
\end{knowledgebox}

\subsection{本章小结}

\begin{itemize}
    \item 注意力机制从2014年的 NMT 应用发展到2017年的 Transformer，再到当今的大语言模型
    \item 自注意力是 Transformer 的核心，允许序列内部各位置相互交互
    \item 降秩乘法注意力直接对应 Transformer 的多头注意力设计
    \item 大语言模型建立在多层 Transformer 之上，通过预训练获得通用能力
\end{itemize}

%% ========== Part 6: NLP 实验实践建议 ==========

